\section{Explicit polylogarithms and ordinary multiple sums}
\label{dhr:sec:explicit}

In this section use the canonical critical weight
$w_r(n)=1/(n+1)_{2r-1}$. Abbreviate $S_{r,N}=S_{w_r,N}$,
$K_{r,N}=S_{r,N}-S_{r,N-1}$, and
\begin{equation}\label{dhr:eq:canonical-GI}
 G_r(z)=\sum_{N\geq0}K_{r,N}z^{N+1},\qquad
 I_{r,q}=\sum_{0\leq n_1<\cdots<n_r}
 \frac{\Delta(\bm n)^2}
 {(n_r+1)^q\prod_i(n_i+1)_{2r-1}},\quad q\geq1.
\end{equation}
Thus $I_{r,q}=\Theta^{-q}G_r(1)$, an ordinary convergent endpoint.
The one-variable case is $G_1=\Li_1$ and $I_{1,q}=\zeta(q+1)$.

\subsection{Two variables: complete cutoff and generating function}
The three moments needed for $r=2$ are
\begin{align}
 M_0(N)&=\frac{(N+1)(N+4)}{4(N+2)(N+3)},\notag\\
 M_1(N)&=\frac{N(N+1)}{4(N+2)(N+3)},\notag\\
 M_2(N)&=H_{N+1}-\frac{11}{4}
                      +\frac1{2(N+2)}+\frac9{2(N+3)}.
 \label{dhr:eq:r2-moments}
\end{align}
They follow directly from the residues of $n^k/(n+1)_3$.
Taking $M_0M_2-M_1^2$ gives the finite identity
\begin{equation}\label{dhr:eq:r2-cutoff}
 \boxed{S_{2,N}=
 \frac{(N+1)(N+4)}{4(N+2)(N+3)}H_{N+1}
 -\frac{(N+1)^2(3N+8)}{4(N+2)^2(N+3)}.}
\end{equation}
This proves the cutoff formula without any asymptotic summation.
It also independently recovers the residue $1/4$ and constant
$(\gamma-3)/4$.

Applying the exact transforms in Section~\ref{dhr:sec:depth} yields
\begin{align}
 G_2(z)={}&\frac{(1-z)^2}{4z^2}\log^2(1-z)
        +\frac{z-1}{2z}\Li_2(z)\notag\\
 &-\frac{6-5z}{4z^2}\log(1-z)+\frac34-\frac3{2z}.
 \label{dhr:eq:r2-G}
\end{align}
The expression is analytic at zero and starts $z^2/144$,
corresponding to the unique first configuration $(0,1)$.
Every negative-power coefficient cancels exactly.
At $z=1/2$, the classical dilogarithm reflection formula
\cite{dlmf:polylog} gives
\begin{equation}\label{dhr:eq:G-half}
 \sum_{0\leq m<n}\frac{(n-m)^2\,2^{-(n+1)}}{(m+1)_3(n+1)_3}
 =\frac12\log^22+\frac72\log2-\frac{\pi^2}{24}-\frac94.
\end{equation}
At the other endpoint, \eqref{dhr:eq:r2-G} gives
\begin{equation}\label{dhr:eq:Abel-constant}
 G_2(z)=-\frac14\log(1-z)-\frac34+o(1),\qquad z\uparrow1.
\end{equation}
The Abel constant is $-3/4$, not the hard-cutoff constant
$(\gamma-3)/4$. This is consistent with
$\sum_{N\geq0}H_{N+1}z^{N+1}=-\log(1-z)/(1-z)$.

\begin{theorem}[An ordinary double-sum evaluation]
\label{dhr:thm:first-evaluation}
\begin{equation}\label{dhr:eq:I21}
 \boxed{\displaystyle
 \sum_{0\leq m<n}
 \frac{(n-m)^2}{(m+1)_3(n+1)_3(n+1)}
 =\zeta(3)-\frac{\pi^2}{24}-\frac34.}
\end{equation}
\end{theorem}
\begin{proof}
We give a telescoping proof in addition to the polylogarithmic
primitive method. Put $v=N+1$. Differencing \eqref{dhr:eq:r2-cutoff}
gives
\[
 K_{2,v-1}=\frac{4(v+1)H_v+v^3-3v^2-6v}
 {4v(v+1)^2(v+2)}.
\]
Partial fractions of $K_{2,v-1}/v$ give a harmonic coefficient
\[
 -\frac3{4v}+\frac1{v+1}-\frac1{4(v+2)}+\frac1{2v^2}
\]
and a nonharmonic coefficient
\[
 -\frac3{4v}+\frac5{4(v+1)}-\frac1{2(v+2)}
                                      +\frac1{2(v+1)^2}.
\]
Group the first three terms of each coefficient into telescoping
differences before summing. Absolute convergence then permits
\begin{align*}
 \sum_{v\geq1}\frac{H_v}{v(v+1)}&=\zeta(2),&
 \sum_{v\geq1}\frac{H_v}{(v+1)(v+2)}&=1,&
 \sum_{v\geq1}\frac{H_v}{v^2}&=2\zeta(3).
\end{align*}
The first two identities follow by summing over the inner index in
$H_v$ and telescoping; the third is
$\zeta(2,1)+\zeta(3)=2\zeta(3)$ by
Lemma~\ref{dhr:lem:Euler}. The harmonic part therefore sums to
$-3\zeta(2)/4+1/4+\zeta(3)$. The nonharmonic part sums to
$\zeta(2)/2-1$. Their total is
$\zeta(3)-\zeta(2)/4-3/4$, proving the theorem.
\end{proof}

Two further primitives give
\begin{align}
 I_{2,2}&=\frac74-\frac{\pi^2}{16}-\frac32\zeta(3)
                            +\frac{\pi^4}{144},\label{dhr:eq:I22}\\
 I_{2,3}&=-\frac52+\frac{17\pi^2}{96}+\zeta(3)
             -\frac{\pi^4}{96}-\frac{\pi^2\zeta(3)}{12}
             +\frac32\zeta(5).
 \label{dhr:eq:I23}
\end{align}
These follow by dividing the same rational-plus-harmonic mass by
$v^2$ and $v^3$, applying \eqref{dhr:eq:T-transform}--\eqref{dhr:eq:U-shift},
and then \eqref{dhr:eq:Euler-zeta}. The delivered exact formula file
contains the full primitive reductions, not just their endpoints.

\subsection{Three variables: a depth-two formula of weight three}
For $r=3$, the determinant gives
\begin{equation}\label{dhr:eq:r3-P}
 P_3(N)=\frac{N(N+1)^2(N^2+17N+54)}
 {10368(N+3)^2(N+4)^2(N+5)}
\end{equation}
and
\begin{equation}\label{dhr:eq:r3-Q}
 Q_3(N)=-\frac{N(N+1)^3(9N^3+125N^2+514N+648)}
 {20736(N+2)(N+3)^3(N+4)^2(N+5)},
\end{equation}
with $S_{3,N}=P_3(N)H_{N+1}+Q_3(N)$. Hence its hard-cutoff
constant is $(\gamma-9/2)/10368$, again agreeing with the independent
Gamma resonance calculation.

For compact notation let $L_j=\Li_j(z)$ and
$D_j=\Li_{j,1}(z,1)$. The generating function is
\begin{align}
 G_3(z)={}&-\frac{5(z-1)^2(7z-4)}{1728z^4}D_1
        +\frac{(z-1)(z+1)}{144z^3}D_2\notag\\
 &-\frac{117z^3+95z^2-513z+300}{10368z^4}L_1\notag\\
 &-\frac{5(z-1)(7z-18)}{1728z^3}L_2
        +\frac{z-1}{144z^2}L_3
        -\frac{281z^2-78z-200}{6912z^3}.
 \label{dhr:eq:r3-G}
\end{align}
Its first nonzero coefficient is $z^3/54432000$, the configuration
$(0,1,2)$. Thus three ordered summation variables have not produced
depth three. The theorem explains this before any coefficient
calculation is performed.

The first three ordinary primitive endpoints are
\begin{align}
 I_{3,1}&=\frac{1903}{124416}
          -\frac{35\pi^2}{62208}-\frac7{864}\zeta(3),
 \label{dhr:eq:I31}\\
 I_{3,2}&=-\frac{6775}{373248}
          +\frac{67\pi^2}{186624}+\frac7{576}\zeta(3),
 \label{dhr:eq:I32}\\
 I_{3,3}&=\frac{100297}{4478976}
          -\frac{403\pi^2}{1492992}-\frac{131}{7776}\zeta(3)
          +\frac{\pi^4}{186624}.
 \label{dhr:eq:I33}
\end{align}
In particular, \eqref{dhr:eq:I31} is the explicit triple sum
\[
 \sum_{0\leq i<j<k}
 \frac{(j-i)^2(k-i)^2(k-j)^2}
 {(i+1)_5(j+1)_5(k+1)_5(k+1)}.
\]
The displayed value is a reduction of this entire sum, not a bound or
an asymptotic approximation.

\subsection{Four variables and the uniform algorithm}
A first four-variable endpoint is
\begin{equation}\label{dhr:eq:I41}
 I_{4,1}=\frac{630863}{268738560000}
 +\frac{14717\pi^2}{67184640000}
 -\frac{487}{124416000}\zeta(3)
 +\frac{\pi^4}{497664000}.
\end{equation}
The definition \eqref{dhr:eq:canonical-GI} specifies its six squared
difference factors and four degree-seven factorial denominators.
The numbers in these expressions can cancel substantially; direct
low-precision evaluation is therefore a poor way to assess them.
The exact rational reduction precedes the high-precision diagnostic.

For any requested $r$ and primitive order $q$, there is a terminating
algorithm: compute the $2r-1$ finite moments by residues, take the two
determinants \eqref{dhr:eq:PQ-det}, difference the cutoff, divide by
$(N+1)^q$, and apply the finite partial-fraction transform. It outputs
both a functional polylogarithmic identity and, when $q\geq1$, its
ordinary zeta endpoint. The implementation in \texttt{code/reductions.py}
uses exact rational arithmetic. It is not restricted in its mathematics
to the ranks and primitive orders printed here, although symbolic
expression growth limits convenient computation at large rank.

The fact that some displayed endpoints omit the highest permitted
weight is an additional cancellation in those cases. We assert the
proved upper bound $r+q$, not a universal sharp-weight formula and
not independence of the surviving zeta coordinates.
